\documentclass[12pt, a4paper, oneside]{ctexart}
\usepackage{amsmath, amsthm, amssymb, bm, xcolor, tikz, framed, graphicx, hyperref, mathrsfs}
\usepackage{bookmark}
\usepackage{pgfplots}
\pgfplotsset{compat=1.18}

\newcommand{\bra}[1]{\langle #1 |}
\newcommand{\ket}[1]{| #1 \rangle}
\newcommand{\braket}[2]{\langle #1 | #2 \rangle}

\title{Quantum Mechanics Homework (Week 13)}
\author{Tang Yi \\ [学号已脱敏]}
\usepackage[a4paper, left=2cm, right=2cm, top=2cm, bottom=2cm]{geometry}
\linespread{1.3} % 行距更紧凑
\definecolor{shadecolor}{RGB}{241, 241, 255}

\newcounter{problemname}
\newenvironment{problem}
  {\begin{shaded}\stepcounter{problemname}\par\noindent\textbf{题目 \theproblemname. }}
  {\end{shaded}\par}
\newenvironment{solution}
  {\par\noindent\textbf{题目 \theproblemname 的解答. }}
  {\par}
\newenvironment{note}
  {\par\noindent\textbf{题目 \theproblemname 的注记. }}
  {\par}

\begin{document}

\maketitle


\begin{problem}
  电子处于二维各向同性谐振子场
\[
V=\frac{1}{2}m\omega_0^2(x^2+y^2)
\]
中，又受到沿 $z$ 轴方向的匀强磁场 $B$ 的作用。

(1) 求电子的能级和本征函数；

(2) 分别讨论 $B\to0$ 和 $B\to\infty$ 两种极限情况以及能级简并度的变化。
\end{problem}

\begin{solution}
  取对称规范
\[
\boldsymbol{A}=\frac{B}{2}(-y,x,0),
\]
电子的哈密顿量为
\[
H=\frac{1}{2m}\left(\boldsymbol{p}+\frac{e}{c}\boldsymbol{A}\right)^2
+\frac{1}{2}m\omega_0^2(x^2+y^2).
\]
整理得
\[
H=\frac{p_x^2+p_y^2}{2m}
+\frac{1}{2}m\Omega^2(x^2+y^2)
-\frac{\omega_c}{2}L_z,
\]
其中
\[
\omega_c=\frac{eB}{mc},\qquad
\Omega=\sqrt{\omega_0^2+\frac{\omega_c^2}{4}}.
\]

能级为
\[
E_{n_+,n_-}=\hbar\Omega(n_++n_-+1)
-\frac{\hbar\omega_c}{2}(n_+-n_-),
\quad n_\pm=0,1,2,\dots
\]
本征函数为二维谐振子本征函数与角动量本征态的组合。

当 $B\to0$ 时，
\[
E=\hbar\omega_0(n_x+n_y+1),
\]
能级简并度为 $n+1$。

当 $B\to\infty$ 时，
\[
E\simeq \hbar\omega_c\left(n_++\frac12\right),
\]
能级退化为朗道能级，简并度随磁场增强而增大。
\end{solution}

\begin{problem}
  考虑 $\sigma_z$ 表象下的自旋状态：
\[
\chi=
\begin{pmatrix}
a\\
b
\end{pmatrix}
\]

(1) 若沿 $z$ 方向进行测量，测得自旋的取值和相应的几率分别为多少？

(2) 若沿 $x$ 方向进行测量，测得自旋的取值和相应的几率分别为多少？

(3) 若先沿 $z$ 方向进行测量，测得自旋向上，即取值为 $\hbar/2$，再沿 $z$ 方向进行测量，测得自旋的取值和相应的几率分别为多少？

(4) 若先沿 $z$ 方向进行测量，测得自旋向上，即取值为 $\hbar/2$，再沿 $x$ 方向进行测量，测得自旋的取值和相应的几率分别为多少？
\end{problem}

\begin{solution}
  (1) 沿 $z$ 方向测量：
\[
P\left(+\frac{\hbar}{2}\right)=|a|^2,\qquad
P\left(-\frac{\hbar}{2}\right)=|b|^2.
\]

(2) 沿 $x$ 方向测量，
\[
|\pm\rangle_x=\frac{1}{\sqrt2}
\begin{pmatrix}1\\ \pm1\end{pmatrix}.
\]
概率为
\[
P\left(\pm\frac{\hbar}{2}\right)
=\frac12|a\pm b|^2.
\]

(3) 先沿 $z$ 方向测量并得到 $+\hbar/2$ 后，态坍缩为
\[
\chi=\begin{pmatrix}1\\0\end{pmatrix},
\]
再次沿 $z$ 测量必得
\[
+\frac{\hbar}{2},\qquad P=1.
\]

(4) 在上述态下沿 $x$ 测量：
\[
P\left(\pm\frac{\hbar}{2}\right)=\frac12.
\]
\end{solution}

\begin{problem}
  考虑泡利算符沿任一方向 $(\theta,\varphi)$ 的分量 $\sigma_n$，在 $\sigma_z$ 表象下为：
\[
\sigma_n=\boldsymbol{\sigma}\cdot\boldsymbol{n}
=\sigma_x n_x+\sigma_y n_y+\sigma_z n_z
=
\begin{pmatrix}
\cos\theta & \sin\theta\,e^{-i\varphi} \\
\sin\theta\,e^{i\varphi} & -\cos\theta
\end{pmatrix}
\]
其中
\[
\boldsymbol{n}=(\sin\theta\cos\varphi,\ \sin\theta\sin\varphi,\ \cos\theta)
\]
是 $(\theta,\varphi)$ 方向的单位矢。

(1) 求其本征值和本征态；

(2) 在 $\sigma_z$ 的本征值为 $1$ 的本征态下，求 $\sigma_n$ 的可能测值、相应的概率和期望值；

(3) 在 $\sigma_n$ 的本征值为 $1$ 的本征态下，求 $\sigma_x,\sigma_y,\sigma_z$ 的可能测值、相应的概率和期望值。
\end{problem}

\begin{solution}
  (1) $\sigma_n$ 的本征值为
\[
\lambda=\pm1.
\]
对应本征态为
\[
|+\rangle_n=
\begin{pmatrix}
\cos\frac{\theta}{2} \\
e^{i\varphi}\sin\frac{\theta}{2}
\end{pmatrix},
\qquad
|-\rangle_n=
\begin{pmatrix}
-\sin\frac{\theta}{2} \\
e^{i\varphi}\cos\frac{\theta}{2}
\end{pmatrix}.
\]

(2) 在 $\sigma_z=1$ 的本征态
\[
|+\rangle_z=\begin{pmatrix}1\\0\end{pmatrix}
\]
下，
\[
P(\sigma_n=\pm1)=\frac12(1\pm\cos\theta),
\]
期望值为
\[
\langle\sigma_n\rangle=\cos\theta.
\]

(3) 在 $\sigma_n=1$ 的本征态下：
\[
\langle\sigma_x\rangle=\sin\theta\cos\varphi,
\quad
\langle\sigma_y\rangle=\sin\theta\sin\varphi,
\quad
\langle\sigma_z\rangle=\cos\theta.
\]
\end{solution}

\begin{problem}
  两个自旋为 $1/2$ 的全同粒子处于二维谐振子势中，记及自旋自由度，写出体系能量最低的三个能级、对应的本征态和简并度。
\end{problem}

\begin{solution}
  单粒子二维谐振子能级为
\[
E=(n_x+n_y+1)\hbar\omega.
\]

基态：
\[
E_0=2\hbar\omega,
\]
空间对称，自旋反对称（单态），简并度 $1$。

第一激发态：
\[
E_1=3\hbar\omega,
\]
空间反对称，自旋对称（三重态），简并度 $3$。

第二激发态：
\[
E_2=4\hbar\omega,
\]
含空间与自旋多种组合，总简并度 $6$。
\end{solution}

\begin{problem}
  两个自旋 $1/2$ 粒子构成的体系，其哈密顿量为
\[
H=J\,\boldsymbol{s}_1\cdot\boldsymbol{s}_2 .
\]
设 $t=0$ 时体系处于状态 $|\uparrow\rangle_1|\downarrow\rangle_2$，求 $t>0$ 时：

(1) 粒子 1 自旋向上的几率；

(2) 粒子 1 和粒子 2 自旋均向上的几率；

(3) 总自旋 $\boldsymbol{S}^2,\,S_z$ 的可能测值和相应的几率；

(4) $\boldsymbol{s}_1,\boldsymbol{s}_2$ 的期望值。
\end{problem}

\begin{solution}
  解：

利用
\[
\boldsymbol{s}_1\cdot\boldsymbol{s}_2
=\frac12(S^2-s_1^2-s_2^2),
\]
三重态能量
\[
E_t=\frac{J}{4},\qquad
\]

单态能量
\[
E_s=-\frac{3J}{4}.
\]

初态分解为
\[
|\uparrow\downarrow\rangle
=\frac{1}{\sqrt2}(|1,0\rangle+|0,0\rangle).
\]

(1)
\[
P_{1\uparrow}(t)=\frac12\left[1+\cos\left(\frac{Jt}{\hbar}\right)\right].
\]

(2)
\[
P_{\uparrow\uparrow}(t)=0.
\]

(3)
\[
S^2=2\hbar^2\ (P=\tfrac12),\quad
S^2=0\ (P=\tfrac12),
\quad
S_z=0\ (P=1).
\]

(4)
\[
\langle\boldsymbol{s}_1\rangle
=\frac{\hbar}{2}\cos\left(\frac{Jt}{\hbar}\right)\hat{z},
\quad
\langle\boldsymbol{s}_2\rangle
=-\langle\boldsymbol{s}_1\rangle.
\]

\end{solution}

\begin{problem}
  证明贝尔基是 $\sigma_{1x}\sigma_{2x}$ 的本征态，并求其本征值。
\end{problem}

\begin{solution}
  四个贝尔态为
\[
|\Phi^\pm\rangle=\frac{1}{\sqrt2}
(|\uparrow\uparrow\rangle\pm|\downarrow\downarrow\rangle),
\]
\[
|\Psi^\pm\rangle=\frac{1}{\sqrt2}
(|\uparrow\downarrow\rangle\pm|\downarrow\uparrow\rangle).
\]

作用算符 $\sigma_{1x}\sigma_{2x}$ 得
\[
\sigma_{1x}\sigma_{2x}|\Phi^\pm\rangle=+|\Phi^\pm\rangle,
\]
\[
\sigma_{1x}\sigma_{2x}|\Psi^\pm\rangle=-|\Psi^\pm\rangle.
\]

故贝尔基是 $\sigma_{1x}\sigma_{2x}$ 的本征态，
本征值分别为 $\pm1$。
\end{solution}
\end{document}